\documentclass[10pt, handout]{beamer}
\beamertemplatenavigationsymbolsempty
\setbeamertemplate{footline}[page number]{}
\usetheme{Madrid}
\usepackage{appendixnumberbeamer}
\usepackage{amssymb}
\usepackage{amsthm}
\usepackage{bbm}
\usepackage[normalem]{ulem}  % for \sout
\newcommand{\ones}{\mathbf{1}}
\newcommand{\indep}{\rotatebox[origin=c]{90}{$\models$}}
\usepackage{qrcode}

%% === basic packages ===
\usepackage{latexsym,multirow}
\usepackage{amssymb,amsmath,bm}
\usepackage{graphicx}
\usepackage[color,all,import,arrow]{xy}
\usepackage{enumerate}
\usepackage{booktabs}
\usepackage{transparent}
\usepackage{tikz}
\usepackage{graphicx}
\usetikzlibrary{calc}


\theoremstyle{plain}
\newtheorem{assumption}{Assumption}
\newtheorem{property}{Property}
\newtheorem{proposition}{Proposition}
\newtheorem{remark}{Remark}
\newtheorem{defi}{Definition}
\newtheorem{thm}{Theorem}




\usepackage{colortbl}  % Required for \rowcolor
\usepackage{xcolor} 

%% == tikz
\usepackage{tikz}
\usepackage{inputenc}
\usetikzlibrary{shapes,arrows,trees,fit,positioning,shapes.misc,calc}

% === some special symbols
\newcommand{\argmax}{\operatornamewithlimits{argmax}}
\newcommand{\argmin}{\operatornamewithlimits{argmin}}
\newcommand{\ind}{\mbox{$\perp\!\!\!\perp$}}
\newcommand{\nind}{\mbox{$\not\perp\!\!\!\perp$}}


\newcommand{\cX}{\mathcal{X}}

\newcommand{\hbx}{\bar{\bm{x}}}
\newcommand{\hbX}{\overline{\bm{X}}}
\newcommand{\hbz}{\bar{\bm{z}}}
\newcommand{\hbZ}{\overline{\bm{Z}}}

\newcommand{\bx}{\bm{x}}
\newcommand{\bX}{\bm{X}}
\newcommand{\ubX}{\underline{\bm{X}}}
\newcommand{\td}{\text{d}}
\newcommand{\bZ}{\bm{Z}}
\newcommand{\ubZ}{\underline{\bm{Z}}}
\newcommand{\bD}{\bm{D}}
\newcommand{\cD}{\mathcal{D}}
\newcommand{\ubD}{\underline{\bm{D}}}
\newcommand{\ucD}{\underline{\mathcal{D}}}
\newcommand{\bz}{\bm{z}}
\newcommand{\ubz}{\underline{\bm{z}}}
\newcommand{\bd}{\bm{d}}
\newcommand{\ubd}{\underline{\bm{d}}}
\newcommand{\bY}{\bm{Y}}
\newcommand{\ubY}{\underline{\bm{Y}}}
\newcommand{\ucY}{\underline{\mathcal{Y}}}
\newcommand{\by}{\bm{y}}
\newcommand{\uby}{\underline{\bm{y}}}
\newcommand{\cY}{\mathcal{Y}}
\newcommand{\add}{\textsc{Add}}
\newcommand{\std}{\textsc{Std}}
\newcommand{\cof}{\textsc{Cof}}

\newcommand{\oD}{\overline{D}}
\newcommand{\oY}{\overline{Y}}

\newcommand{\bS}{\bar{S}}
\newcommand{\bs}{\bar{s}}
\newcommand{\Ss}{\mathcal{S}}
\newcommand{\N}{\mathbb{N}}
\newcommand{\E}{\mathbb{E}}
\newcommand{\R}{\mathbb{R}}
\newcommand{\bI}{\bm{I}}
\newcommand{\bH}{\bm{H}}
\newcommand{\bM}{\bm{M}}
\newcommand{\bP}{\bm{P}}
\newcommand{\bQ}{\bm{Q}}
\newcommand{\bV}{\bm{V}}
\newcommand{\bU}{\bm{U}}
\newcommand{\bW}{\bm{W}}
\newcommand{\bw}{\bm{w}}
\def\independenT#1#2{\mathrel{\rlap{$#1#2$}\mkern2mu{#1#2}}}
\DeclareMathOperator{\sgn}{sgn}
\DeclareMathOperator{\diag}{diag}
\providecommand{\norm}[1]{\lVert#1\rVert}


\newcommand{\cU}{\mathcal{U}}

\newcommand{\bepsilon}{\bm{\epsilon}}
\newcommand{\boldeta}{\bm{\eta}}
\newcommand{\bC}{\bm{C}}
\newcommand{\bl}{\bm{\ell}}
\newcommand{\ATT}{\mathrm{ATT}}
\newcommand{\Yone}{Y(1)}
\newcommand{\Yzero}{Y(0)}
\newcommand{\ATE}{\mathrm{ATE}}
\usepackage{booktabs,tabularx}
\newcommand{\Var}{\mathrm{Var}}
\newcommand{\Cov}{\mathrm{Cov}}


% BK: Added Userpackages and Commands
\usepackage{dsfont}
\usepackage{caption}
\usepackage{tikz-cd}
\usetikzlibrary{arrows.meta}

\newcommand{\lt}{\left}
\newcommand{\rt}{\right}
\newcommand{\indicator}[1]{\mathds{1}{\left\{ #1 \right\}}}
\newcommand{\prb}[1]{\mathbb{P}{\left( #1 \right)}}
\newcommand{\prbcond}[2]{\mathbb{P}{\left( #1 \mid #2 \right)}}
\newcommand{\Eb}[1]{\mathbb{E}{\left[ #1 \right]}}
\newcommand{\Ebcond}[2]{\mathbb{E}{\left[ #1 \mid #2 \right]}}
\newcommand{\indist}{\overset{d}{\to}}
\newcommand{\inprob}{\overset{p}{\to}}
\newcommand{\as}{\overset{a.s.}{\to}}
\newcommand{\Nb}[1]{\mathcal{N}{\left( #1 \right)}}
\newcommand{\abs}[1]{\left|#1\right|}
\newcommand{\cF}{\mathcal{F}}

\usepackage[style=authoryear]{biblatex}

% Add bib resource
\addbibresource{references.bib}


\usepackage{booktabs}
\usepackage[scale=2]{ccicons}

\usepackage{pgfplots}
\usepgfplotslibrary{dateplot}

\usepackage{xspace}
\newcommand{\themename}{\textbf{\textsc{metropolis}}\xspace}


\title{Midterm Review Lecture STAT 286}
\author{Benedikt Koch}
\date{October 8th, 2025}
\begin{document}

\begin{frame}[plain]

  \maketitle

\end{frame}


\AtBeginSection[]
{
  \begin{frame}{Outline}
    \tableofcontents[currentsection]
  \end{frame}
}


\section{Introduction}

\begin{frame}{Causal Inference Pipeline}
\begin{itemize}
    \item \textbf{Setup}: Binary treatment \(T_i\in\{0,1\}\); outcome \(Y_i\) may be binary, count, or continuous.
    \pause
    \item \textbf{Goal:} Quantify the causal effect of \(T\) on \(Y\).
    \pause
    \item \textbf{Estimand}: 
    
    \(\displaystyle \tau_{\text{SATE}}=\frac{1}{n}\sum_{i=1}^n\big(Y_i(1)-Y_i(0)\big), \qquad  \tau_{\text{PATE}}=\mathbb{E}\big[Y(1)-Y(0)\big]\)
    \pause
    \item \textbf{Identification:} A causal estimand \(\tau\) is \emph{identified} if there exists a known map \(g\) such that
  \[
    \tau=g(\mathbb{P}),\qquad \mathbb{P} \text{ is the observable law of } (Y,T,X).
  \]
  Requires assumptions!
    \item \textbf{Example (under CRD):}
  \[
    \mathbb{E}[Y(1)-Y(0)] \;=\; \mathbb{E}[Y\!\mid\!T{=}1]-\mathbb{E}[Y\!\mid\!T{=}0].
  \]
  \pause
    \item \textbf{Estimation \& inference:} Estimate the identified estimand (permutation + test inversion, diff-in-means + Neyamnn variance, OLS\,+\,HC2), report s.e./CI, hope its significant
\end{itemize}

\end{frame}


\begin{frame}{Assumptions}


\begin{itemize}
  \item \textbf{Causal ordering.} $T_i \to Y_i$
  \pause

  \vspace{1em}
  \item \textbf{Stable Unit Treatment Value Assumption (SUTVA).}
  \setlength{\itemsep}{6pt} % slightly smaller spacing for nested list
    \begin{itemize}
      \item \emph{No interference:} $Y_i(T_1,\dots,T_N)=Y_i(T_i)$
      \item \emph{Consistency:} $Y_i = Y_i(T_i)$
    \end{itemize}

  \pause
  \vspace{1em}
  \item \textbf{Randomization.}
  \setlength{\itemsep}{6pt}
    \begin{itemize}
        \item \emph{Unconfoundedness:} $(Y_i(1),Y_i(0)) \;\perp\!\!\!\perp\; T_i$
        \item \emph{Overlap:} $0<\Pr(T_i = 1) <1$, in CRD: $\Pr(T_i=1) = \frac{n_1}{n}$
    \end{itemize}
\end{itemize}
\end{frame}


\begin{frame}{When Assumptions Fail}
\textbf{Scenario:} Effect of a flu shot ($T$) on days sick ($Y$).

\pause
\vspace{1em}
\begin{itemize}
  \item \textbf{Causal ordering}: People vaccinate \emph{after} feeling ill $\Rightarrow Y \to T$.
  
  \pause
  \vspace{1em}
  \item \textbf{No interference}: Herd immunity; $Y_i$ depends on others’ $T_j$.

  \pause
  \vspace{1em}
  \item \textbf{Consistency}: “Flu shot” has different versions $\Rightarrow T{=}1$ not well-defined.

  \pause
  \vspace{1em}
  \item \textbf{Unconfoundedness}: Health-consciousness or access affects both $T$ and $Y$ (selection bias).

  \pause
  \vspace{1em}
  \item \textbf{Overlap}: Some strata always/never vaccinate; no comparable counterfactuals.
\end{itemize}

\medskip
\textbf{Result:} $\mathbb{E}[Y(1)-Y(0)] \neq \E[Y\mid T{=}1]-\E[Y\mid T{=}0]$
\end{frame}

\section{Permutation Test}

\begin{frame}{Permutation Test (1/4)}

\begin{itemize}
  \item \textbf{Setup}:
  \setlength{\itemsep}{6pt} % slightly tighter spacing for nested list
    \begin{itemize}
      \item Units: $i = 1,\ldots,n$
      \item Treatment: $T_i \in \{0,1\}$
      \item Outcome: $Y_i = Y_i(T_i)$, \quad $\mathcal{O}_n=\{Y_i(0),Y_i(1)\}_{i=1}^n$
      \item Completely Randomized Design (CRD): $\Pr(T_i=1\mid \mathcal{O}_n)=\Pr(T_1=1)=\tfrac{n_1}{n}$
    \end{itemize}

  \pause
  \vspace{1em}
  \item \textbf{Sharp Null}: $H_0:\; Y_i(1)-Y_i(0)=\tau_0$ \; (often $\tau_0=0$)

  \pause
  \vspace{1em}
  \item \textbf{Test Statistic}:
  $S_{\tau_0}=\frac{1}{n_1}\sum_{i=1}^n T_i\,Y_i(T_i)\;-\;\frac{1}{n_0}\sum_{i=1}^n (1-T_i)\,(Y_i(T_i) + \tau_0)$ \quad (many possibilities, calibrate $\E[S_{\tau_0} \mid H_0] = 0$ to make inference easier (not necessary))

  \pause
  \vspace{1em}
  \item \textbf{Inference}: Permutation with inversion
  
\end{itemize}
\end{frame}

\begin{frame}{Permutation Test (2/4)}

\small
\begin{itemize}
  \item \textbf{Sharp Null:}\; $H_0:\; Y_i(1)-Y_i(0)=\tau_0$
  \vspace{0.5em}
  \item \textbf{Test Statistic:}\;
  \[
  S_{\tau_0} = \frac{1}{n_1}\sum_{i=1}^n T_i\,Y_i(T_i)
  \;-\;
  \frac{1}{n_0}\sum_{i=1}^n (1-T_i)\,(Y_i(T_i) + \tau_0)
  \]
\end{itemize}

\pause
\vspace{0.8em}
\textbf{Pipeline.}
\begin{itemize}
  \setlength{\itemsep}{6pt}
  \item Draw $b = 1,\ldots,B$ random assignments $\bm{T}^{(b)} \in \{0,1\}^n$.
  \item Compute potential outcomes under each assignment:
  \begin{itemize}
    \setlength{\itemsep}{4pt}
    \item If $T_i^{(b)} = T_i$: $Y_i(T_i^{(b)}) = Y_i(T_i)$
    \item If $(T_i^{(b)}, T_i) = (0,1)$: $Y_i(T_i^{(b)}) = Y_i(0) = Y_i(1) - \tau_0 = Y_i(T_i) - \tau_0$
    \item If $(T_i^{(b)}, T_i) = (1,0)$: $Y_i(T_i^{(b)}) = Y_i(1) = Y_i(0) + \tau_0 = Y_i(T_i) + \tau_0$
  \end{itemize}
  \item Recalculate the test statistic:
  \[
  S^{(b)} = 
  \frac{1}{n_1}\sum_{i=1}^n T_i^{(b)}\,Y_i(T_i^{(b)})
  - 
  \frac{1}{n_0}\sum_{i=1}^n (1-T_i^{(b)})\,(Y_i(T_i^{(b)}) + \tau_0).
  \]
\end{itemize}
\end{frame}

\begin{frame}{Permutation Test (3/4)}

\footnotesize
\textbf{p-values.}
\begin{itemize}
  \setlength{\itemsep}{5pt}
  \item $H_1:\ \tau \neq \tau_0$ 
  \[
  p_{\text{perm}}(\tau_0)
  = \frac{1+\sum_{b=1}^B \indicator{|S^{(b)}-\mu|\ge|S_{\tau_0}^{\text{obs}}-\mu|}}{B+1}
  \]
  ($\mu = \E[S_{\tau_0} \mid H_0(\tau_0)]$; here $\mu = 0$; alternative $p_{\text{perm}}(\tau_0) = 2 \cdot \min(p_{\text{perm},+}(\tau_0), p_{\text{perm},-}(\tau_0))$)
  \item $H_1:\ \tau>\tau_0$ 
  \[
  p_{\text{perm},+}(\tau_0)
  = \frac{1+\sum_{b=1}^B \indicator{S^{(b)}\ge S_{\tau_0}^{\text{obs}}}}{B+1}.
  \]
  \item $H_1:\ \tau<\tau_0$
  \[
  p_{\text{perm},-}(\tau_0)
  = \frac{1+\sum_{b=1}^B \indicator{S^{(b)}\le S_{\tau_0}^{\text{obs}}}}{B+1}.
  \]
\end{itemize}

\pause

\textbf{Confidence set.}
\[
\mathcal C_{1-\alpha}=\{\tau:\ p(\tau)>\alpha\},
\]
where \(p(\tau)\) is the p\textendash value for \(H_0(\tau)\).
If the test statistic is monotone in \(\tau\), then \(\mathcal C_{1-\alpha}\) is an interval.

\pause

\bigskip
\textbf{Point estimate.}
\[
\hat\tau=\operatorname*{arg\,max}_{\tau}\ p(\tau).
\]

\end{frame}



\begin{frame}{Permutation Test (4/4)}
\begin{itemize}
    \item \textbf{Pros}: $Y_i$ can be binary, categorical, or continuous

    \vspace{0.5em}
    \item \textbf{Cons}: May be less powerful; sensitive to outliers; rank-based versions mitigate this (although they require non-binary $Y_i$)
\end{itemize}
\end{frame}

\begin{frame}{Rank-Sum Test (1/2)}

\begin{itemize}
  \setlength{\itemsep}{6pt}
  \item \textbf{Setup:} \(Y_i\) non-binary (typically continuous).
  \item \textbf{Sharp null:} \(H_0:\; Y_i(1)-Y_i(0)=\tau_0\).
  \item \textbf{Test statistic:} \(S_{\tau_0}=\sum_{i=1}^n T_i\,R_i(\bm Y-\tau\bm T)\).
\end{itemize}

\pause
\vspace{0.8em}
\textbf{Invariance under \(H_0\):} Let \(\bm T,\bm A\in\{0,1\}^n\),
  \[
    \bm Y(\bm T)-\tau_0\bm T \;=\; \bm Y(0) \;=\; \bm Y(\bm A)-\tau_0\bm A.
  \]
  \vspace{-1em}
  \begin{itemize}
    \setlength{\itemsep}{4pt}
    \item \(T_i=1:\; Y_i(T_i)-\tau_0 T_i=Y_i(1)-\tau_0=Y_i(0)\).
    \item \(T_i=0:\; Y_i(T_i)-\tau_0 T_i=Y_i(0)\).
  \end{itemize}
  Hence \(S_{\tau_0}\) does not depend on \(\tau_0\) given \(H_0\); use a normal approximation or permutation inference.

\pause
\vspace{0.8em}
\textbf{Permutation Inference}: Draw $b = 1,\ldots,B$ random assignments $\bm{T}^{(b)} \in \{0,1\}^n$.
  \[
  S^{(b)} = \sum_{i=1}^n T^b_i \,R_i(\bm Y-\tau\bm T)
  \]
  Note that we do not change the expression inside the ranks.
\end{frame}

\begin{frame}{Rank-Sum Test (2/2)}

\vspace{0.8em}
\textbf{Other important facts}

\begin{itemize}
  \setlength{\itemsep}{8pt}
    \pause
  \item \emph{Monotonicity:} \(S_\tau\) is nonincreasing in \(\tau\); if \(\tau_1<\tau_2\) then \(S_{\tau_1}>S_{\tau_2}\).

    \pause
  \item \emph{Maximal treatment effects:} \(H^{\max}_0(\tau_0):\ \max_i\{Y_i(1)-Y_i(0)\}\le\tau_0\).
  \[
    \Pr\!\big(S_{\tau_0}\!\ge s \mid H^{\max}_0(\tau_0)\big)
    \;\le\;
    \Pr\!\big(S_{\tau_0}\!\ge s \mid H_0(\tau_0)\big).
  \]
  Same one-sided sharp-null statistic yields valid (conservative) p-values.
\end{itemize}

\end{frame}



% \begin{frame}{Rank-Sum Test}
% \begin{itemize}
%   \item \textbf{Setup:} \(Y_i\) non-binary (typically continuous).
%   \item \textbf{Sharp null:} \(H_0:\; Y_i(1)-Y_i(0)=\tau_0\).
%   \item \textbf{Test statistic:} \(S_\tau=\sum_{i=1}^n T_i\,R_i(\bm Y-\tau\bm T)\).
% \end{itemize}

% \medskip
% \textbf{Important facts}
% \begin{itemize}
%   \item \emph{Invariance under \(H_0\):} for any \(\bm T,\bm A\in\{0,1\}^n\),
%   \[
%     \bm Y(\bm T)-\tau_0\bm T=\bm Y(0)=\bm Y(\bm A)-\tau_0\bm A.
%   \]
%   \vspace{-1.5em}
%   \begin{itemize}
%     \item \(T_i=1:\; Y_i(T_i)-\tau T_i=Y_i(1)-\tau=Y_i(0)\).
%     \item \(T_i=0:\; Y_i(T_i)-\tau T_i=Y_i(0)\).
%   \end{itemize}
%   Hence \(S_{\tau_0}\) does not depend on $\tau_0$ given $H_0$: use a normal approximation or permutation inference.
%   \item \emph{Monotonicity:} \(S_\tau\) is non-increasing in \(\tau\); if \(\tau_1<\tau_2\) then \(S_{\tau_1}>S_{\tau_2}\).
%   \item \emph{Maximal Treatment Effects}: $H^{\max}_0(\tau_0):\ \max_{1\le i\le n}\{Y_i(1)-Y_i(0)\}\ \le\ \tau_0$.
%   \[
%     \Pr(S_{\tau_0} \geq s \mid H^{\max}_0(\tau_0)) \leq \Pr(S_{\tau_0} \geq s \mid H_0(\tau_0)).
%     \]
%     Same test statistic under the one-sided sharp null $H_1: \tau > \tau_0$, gives valid p-values.
% \end{itemize}
% \end{frame}

% \begin{frame}{Inverting the Test}
% \textbf{Confidence interval.}
% \[
% \mathcal C_{1-\alpha}=\{\tau:\ p(\tau)>\alpha\},
% \]
% where \(p(\tau)\) is the p\textendash value for \(H_0(\tau)\).
% If the test statistic is monotone in \(\tau\) (e.g., \(S_\tau\)), then \(\mathcal C_{1-\alpha}\) is an interval.

% \bigskip
% \textbf{Point estimate.}
% \[
% \hat\tau=\operatorname*{arg\,max}_{\tau}\ p(\tau).
% \]
% \end{frame}

\section{Average Treatment Effects}

% --- Slide 1: rows through Estimand ---
\begin{frame}{Finite vs.\ Super - Population (1/3)}
\footnotesize
\begin{table}
\centering
\begin{tabular}{|>{\centering\arraybackslash}p{0.24\textwidth}|>{\centering\arraybackslash}p{0.33\textwidth}|>{\centering\arraybackslash}p{0.33\textwidth}|}
\hline
 & \textbf{Finite-Population} & \textbf{Super-Population} \\
\hline
\textit{Population} &
Set of units in the study is fixed and finite &
Units are drawn i.i.d.\ from a large or infinite ``super-population.'' \\
\hline
\textit{Potential Outcomes} &
Fixed set $\mathcal{O}_n=\{Y_i(0),Y_i(1)\}_{i=1}^n$ &
$(Y_i(0),Y_i(1)) \text{ i.i.d.\ }\sim \mathbb{P}^{\ast}$ \\
\hline
\textit{Inference} &
Treatment effects for this specific group of units &
Treatment effects for the underlying population distribution \\
\hline
\textit{Source of Randomness} &
Treatment assignment &
\begin{tabular}{@{}c@{}}Treatment assignment \\ Sampling of units from the population\end{tabular} \\
\hline
\textit{Random Variables} &
Treatment $T_i$ &
Treatment and Potential Outcomes $T_i, Y_i(0), Y_i(1)$ \\
\hline
\textit{Example} &
What is the ATE on these $n$ units in our experiment? &
What is the ATE in the population from which these $n$ units were sampled? \\
\hline
\textit{Validity} &
Internal &
External \\
\hline
\end{tabular}
\end{table}
\end{frame}

% --- Slide 2: Estimator through Estimator of Variance + Notation ---
\begin{frame}{Finite vs.\ Super - Population (2/3)}
\footnotesize
\begin{table}
\centering
\begin{tabular}{|>{\centering\arraybackslash}p{0.24\textwidth}|>{\centering\arraybackslash}p{0.33\textwidth}|>{\centering\arraybackslash}p{0.33\textwidth}|}
\hline
 & \textbf{Finite-Population} & \textbf{Super-Population} \\
 \hline
\textit{Estimand} &
$\displaystyle \tau_{\mathrm{SATE}}=\frac{1}{n}\sum_{i=1}^n\!\big(Y_i(1)-Y_i(0)\big)$ &
$\displaystyle \tau_{\mathrm{PATE}}=\mathbb{E}\!\left[Y(1)-Y(0)\right]$ \\
\hline
\textit{Estimator} & \multicolumn{2}{c|}{%
\begin{tabular}{c}
Difference in means \\
$\hat\tau = \overline Y_1 - \overline Y_0$
\end{tabular}} \\
\hline
\textit{Unbiasedness} &
$\E[\hat\tau \mid \mathcal{O}_n] = \tau_{\mathrm{SATE}}$ &
$\E[\hat\tau] = \tau_{\mathrm{PATE}}$ \\
\hline
\textit{Variance of Estimator} &
$\displaystyle \mathrm{Var}(\hat\tau\mid \mathcal O_n)
= \frac{S_1^2}{n_1}+\frac{S_0^2}{n_0}-\frac{S_\ast^2}{n} \leq \frac{S_1^2}{n_1}+\frac{S_0^2}{n_0}$
&
$\displaystyle \mathrm{Var}(\hat\tau)=\frac{\sigma_1^2}{n_1}+\frac{\sigma_0^2}{n_0}$ \\
\hline
\textit{Estimator of Variance} &
\multicolumn{2}{c|}{%
\begin{tabular}{c}
$\displaystyle \widehat{\mathrm{Var}}(\hat\tau)=\frac{\hat\sigma_1^2}{n_1}+\frac{\hat\sigma_0^2}{n_0}$
\\
Conservative for FP, Unbiased for SP
\end{tabular}} \\
\hline
\textit{CI in practice} & \multicolumn{2}{c|}{%
\begin{tabular}{c}
$\displaystyle \text{s.e.}=\sqrt{\frac{\hat\sigma_1^2}{n_1}+\frac{\hat\sigma_0^2}{n_0}}$\\
$\displaystyle (1-\alpha)\times 100\%\ \text{CI: }\big[\hat\tau \pm z_{1-\alpha/2}\cdot \text{s.e.}\big]$
\end{tabular}} \\
\hline
\end{tabular}
\end{table}
\end{frame}


\begin{frame}{Finite vs.\ Super - Population (3/3)}
\textbf{Notation}
\[
\begin{aligned}
&n_1=\sum_{i=1}^n T_i,\quad n_0=n-n_1,\qquad
\overline{Y(t)} =\frac{1}{n}\sum_{i=1}^n Y_i(t), \quad \overline Y_t=\frac{1}{n_t}\sum_{i:T_i=t}Y_i, \\
&S_t^2=\frac{1}{n-1}\sum_{i=1}^n\big(Y_i(t)-\overline{Y(t)}\big)^2,\quad 
S_{01}=\frac{1}{n-1}\sum_{i=1}^n\!\big(Y_i(1)-\overline{Y(1)}\big)\big(Y_i(0)-\overline{Y(0)}\big),\\
&\sigma_t^2=\mathrm{Var}\!\big(Y_i(t)\big) = \E[S_t^2],\qquad \hat\sigma_t^2=\frac{1}{n_t-1}\sum_{i:T_i=t}\big(Y_i-\overline Y_t\big)^2.
\end{aligned}
\]
\end{frame}

\section{Linear Regression}

\begin{frame}{Potential Outcomes as Linear Models}

\textbf{Setup.} Binary treatment $T_i\in\{0,1\}$. Define
\[
\alpha_i := Y_i(0), \qquad 
\beta_i := Y_i(1)-Y_i(0), \qquad
\alpha := \E[Y_i(0)], \qquad 
\beta := \E[Y_i(1)-Y_i(0)].
\]

\medskip
\textbf{Decomposition:}
\begin{align*}
\onslide<+->{Y_i(t) 
&= Y_i(0) + \big(Y_i(1)-Y_i(0)\big)\,t }\\[2pt]
\onslide<+->{&= \E[Y_i(0)] + \big(Y_i(0) - \E[Y_i(0)]\big)}\\[2pt]
\onslide<+->{&\quad + \E[Y_i(1)-Y_i(0)]\,t + \big(Y_i(1)-Y_i(0) - \E[Y_i(1)-Y_i(0)]\big)\,t }\\[2pt]
\onslide<+->{&= \alpha + (\alpha_i-\alpha) + \beta\,t + (\beta_i-\beta)\,t }\\[2pt]
\onslide<+->{&= \alpha + \beta\,t + \underbrace{(\beta_i-\beta)\,t + (\alpha_i-\alpha)}_{=:~\varepsilon_i(t)} }\\[2pt]
\onslide<+->{&= \alpha + \beta\,t + \varepsilon_i(t).}
\end{align*}

\medskip
\onslide<+->{\textbf{Key facts.} $\E[\varepsilon_i(t)]=0$. Potential outcomes can be written as linear models without assumptions. Regressing $Y_i\sim[1,\,T_i]$ targets $\beta=\E[Y_i(1)-Y_i(0)]$. Extends to \emph{categorical} $T$. For \emph{continuous} $T$, writing $Y_i(t)=\alpha_i+\beta_i t$ imposes linearity in $t$.}
\end{frame}

\begin{frame}{Review: Linear Regression (1/2)}
\small
\begin{itemize}
  \item \textbf{Linear model:}
  \[
    \bY = \bX \bm{\beta}^\ast + \bm{\varepsilon}.
  \]
  where
  \[
  \bY =
  \begin{bmatrix}
  Y_1 \\ \vdots \\ Y_n
  \end{bmatrix}, \quad
  \bX =
  \begin{bmatrix}
  1 & X_{11} & \cdots & X_{K1} \\
  \vdots & \vdots & \ddots & \vdots \\
  1 & X_{n1} & \cdots & X_{Kn}
  \end{bmatrix}, \quad
  \bm{\beta}^\ast =
  \begin{bmatrix}
  \beta^\ast_0 \\ \vdots \\ \beta^\ast_K
  \end{bmatrix}.
  \]
  \pause
  \vspace{0.5em}
  Note that since we have not specified $\bm{\varepsilon}$, this model is fully general.
  \pause
  \vspace{1em}

  \item \textbf{Ordinary Least Squares (OLS) estimator:}
  \[
    \hat{\bm{\beta}} = \argmin_{\bm{\beta}} \|\bY - \bX \bm{\beta}\|_2,
  \]
  which leads to the closed-form solution:
  \[
    \hat{\bm{\beta}} = (\bX^\top \bX)^{-1} \bX^\top \bY
    = \left(\frac{1}{n}\sum_{i=1}^n \bX_i \bX_i^\top\right)^{-1}
      \left(\frac{1}{n}\sum_{i=1}^n \bX_i Y_i\right),
  \]
  where $\bX_i = [1, X_{1i}, \ldots, X_{Ki}]^\top$.
  \pause
  \vspace{0.5em}
  \textbf{Interpretation:} The OLS estimator gives the \emph{best linear approximation} of $\E[Y_i \mid \bX_i]$.
\end{itemize}
\end{frame}

\begin{frame}{Review: Linear Regression (2/2)}
\small
\textbf{Assumption: Strict exogeneity}
\[
\E[\varepsilon_i \mid \bX_i] = 0.
\]
Note: Under CRD with $\bX_i = [1, T_i]$ this holds. If covariates are included, this does not hold necessarily, but imposes a linearity assumption (see pooled case later on).
\pause
\vspace{0.8em}

\begin{itemize}
  \item \textbf{Unbiasedness:}
  \[
    \E[\hat{\bm{\beta}} \mid \bX] = \bm{\beta}^\ast.
  \]
  \pause
  \vspace{0.8em}

  \item \textbf{Finite-sample variance (sandwich form):}
  \[
    \Var(\hat{\bm{\beta}} \mid \bX)
    = \underbrace{\color{blue}{(\bX^\top \bX)^{-1}}}_{\text{Bread}}
      \;\underbrace{\color{red}{\bX^\top \bm{\Omega} \bX}}_{\text{Meat}}\;
      \underbrace{\color{blue}{(\bX^\top \bX)^{-1}}}_{\text{Bread}},
  \]
  where
  \[
    \bm{\Omega} = \diag(\sigma_1^2, \ldots, \sigma_n^2),
    \qquad \sigma_i^2 = \Var(\varepsilon_i \mid \bX_i).
  \]
  
  \pause
  \vspace{0.5em}
  This is the foundation for all robust (EHW/HC2) and homoskedastic OLS variance formulas.
\end{itemize}
\end{frame}

\begin{frame}{Homoskedasticity}
\small
\begin{itemize}
  \item \textbf{Assumption:} Errors have constant variance  
  \[
  \Var(\varepsilon_i \mid \bX_i) = \sigma^2.
  \]
  \pause
  Then
  \[
  \Var(\hat{\bm{\beta}} \mid \bX)
  = \color{blue}{(\bX^\top \bX)^{-1}}
    \color{red}{\bX^\top \bm{\Omega} \bX}
    \color{blue}{(\bX^\top \bX)^{-1}}
  = \color{red}{\sigma^2}\,\color{blue}{(\bX^\top \bX)^{-1}}.
  \]
  \pause
  \vspace{0.4em}
  This assumption is often unrealistic in applied settings.
  \pause
  \vspace{1em}

  \item \textbf{Example} (binary treatment, CRD): Let $\bX_i = [1, T_i]^\top$.  
  Then
  \[
  (\bX^\top \bX)^{-1}
  = 
  \left(
  \begin{bmatrix}
  \sum_i 1 & \sum_i T_i \\[3pt]
  \sum_i T_i & \sum_i T_i^2
  \end{bmatrix}
  \right)^{-1}
  = \frac{1}{n_1 n_0}
  \begin{bmatrix}
  n_1 & -n_1 \\[3pt]
  -n_1 & n
  \end{bmatrix}.
  \]
  \pause
  Hence
  \[
  \Var(\hat{\beta}_T \mid \bX)
  = \sigma^2\!\left(\frac{1}{n_1} + \frac{1}{n_0}\right).
  \]
  \pause
  \vspace{0.3em}
  Only valid if treatment and control variances are equal:
  \[
  \Var(Y(1)) = \Var(Y(0)) = \sigma^2.
  \]
\end{itemize}
\end{frame}

\begin{frame}{Heteroskedasticity: Eicker–Huber–White (EHW, HC0)}
\small
\begin{itemize}
  \item \textbf{Motivation:} Allow error variances to differ across $i$.
  \pause
  \vspace{0.5em}
  \item \textbf{HC0 (EHW) estimator:}
  \[
  \widehat{\Var}(\hat{\bm{\beta}} \mid \bX)
  = \color{blue}{(\bX^\top \bX)^{-1}}
    \color{red}{\bX^\top \diag(\hat{\varepsilon}_i^2)\,\bX}
    \color{blue}{(\bX^\top \bX)^{-1}},
  \]
  where $\hat{\varepsilon}_i = Y_i - \bX_i^\top \hat{\bm{\beta}}$.
  \pause
  \vspace{1em}

  \item \textbf{Example} (binary treatment, CRD)
  \[
  \Var(\hat{\beta}_T \mid \bX)
  = \frac{\tilde{\sigma}_1^2}{n_1} + \frac{\tilde{\sigma}_0^2}{n_0},
  \]
  where
  \[
  \tilde{\sigma}_t^2 = \frac{1}{n_t}\sum_{i:T_i=t}\big(Y_i-\overline{Y}_t\big)^2.
  \]
  \pause
  \vspace{0.5em}
  \item This estimator is typically too small (biased downward).  
  \pause
  \vspace{0.3em}
  Following Neyman’s correction, we would like to have
  \[
  \hat{\sigma}_t^2
  = \frac{1}{n_t-1}\sum_{i:T_i=t}\big(Y_i-\overline{Y}_t\big)^2.
  \]
  \pause
  \vspace{0.3em}
  HC0 is \emph{consistent} but biased in finite samples.
\end{itemize}
\end{frame}

\begin{frame}{Heteroskedasticity: HC2}
\small
\begin{itemize}
  \item \textbf{Idea:} Correct HC0’s downward bias using leverage adjustment.
  \pause
  \vspace{0.4em}
  \item \textbf{HC2 estimator:}
  \[
  \widehat{\Var}(\hat{\bm{\beta}} \mid \bX)
  = \color{blue}{(\bX^\top \bX)^{-1}}
    \color{red}{\bX^\top \diag\!\left(\frac{\hat{\varepsilon}_i^2}{1 - p_{ii}}\right)\bX}
    \color{blue}{(\bX^\top \bX)^{-1}},
  \]
  where $p_{ii} = \bX_i^\top(\bX^\top\bX)^{-1}\bX_i$ is the leverage of observation $i$.
  \pause
  \vspace{1em}

  \item \textbf{In a CRD with $\bX_i = [1,T_i]^\top$:}
  \[
  p_{ii}
  = \frac{1}{n_1n_0}
    \begin{bmatrix}1\\T_i\end{bmatrix}^\top
    \begin{bmatrix}
      n_1 & -n_1 \\[3pt]
      -n_1 & n
    \end{bmatrix}
    \begin{bmatrix}1\\T_i\end{bmatrix}
  = \frac{1}{n_1}T_i + \frac{1}{n_0}(1-T_i).
  \]
  \pause
  Thus
  \[
  \frac{1}{1-p_{ii}} =
  \begin{cases}
    \frac{n_1}{n_1-1}, & T_i=1,\\[4pt]
    \frac{n_0}{n_0-1}, & T_i=0.
  \end{cases}
  \]
  \pause
  \vspace{0.5em}
  This correction yields the \textbf{Neyman variance}:
  \[
  \Var(\hat{\beta}_T \mid \bX)
  = \frac{\hat{\sigma}_1^2}{n_1} + \frac{\hat{\sigma}_0^2}{n_0},
  \quad
  \hat{\sigma}_t^2 = \frac{1}{n_t-1}\sum_{i:T_i=t}\big(Y_i-\overline{Y}_t\big)^2.
  \]
  \pause
  In very small samples, adjust further by scaling with $\tfrac{n}{n-(K+1)}$ or use \textbf{HC3}.
\end{itemize}
\end{frame}

\begin{frame}{Linear Regression with No Covariates}
\small
\textbf{Setup.} Recall \(Y_i(t) = \alpha + \beta t + \varepsilon_i(t)\).  
Regressing \(Y \sim [1, T]\) yields (without assumptions)
\pause
\begin{align*}
\hat\alpha &= \frac{1}{n_0}\sum_{i=1}^n (1-T_i)Y_i \;=\; \overline{Y}_0, \\[4pt]
\hat\beta_{\text{DiM}} &= \frac{1}{n_1}\sum_{i=1}^n T_i Y_i 
\;-\; \frac{1}{n_0}\sum_{i=1}^n (1-T_i)Y_i.
\end{align*}
\pause
\textbf{Interpretation.}  
\(\hat\beta_{\text{DiM}}\) is the difference-in-means (DiM) estimator.
\pause
\medskip

\textbf{Variance.}  
Using HC2 variance estimation (Neyman’s formula),
\[
\widehat{\Var}(\hat\beta_{\text{DiM}}) = 
\frac{\hat\sigma_1^2}{n_1} + \frac{\hat\sigma_0^2}{n_0}, 
\quad 
\hat\sigma_t^2 = \frac{1}{n_t-1}\sum_{i:T_i=t}(Y_i-\overline{Y}_t)^2.
\]
\pause
Thus, OLS with an intercept and binary \(T_i\) reproduces the classical difference-in-means (DiM) estimator.
\end{frame}

\begin{frame}{Linear Regression with Covariates (Pooled) (1/3)}
\textbf{Motivation.} Even under randomization, covariates can be imbalanced; adjusting for \(X\) can improve precision.

\pause
\vspace{0.6em}
\textbf{Pooled (ANCOVA) model.}
\[
Y_i \;=\; \alpha \;+\; \beta\,T_i \;+\; \bm{\gamma}^\top \bX_i \;+\; \varepsilon_i.
\]

\pause
\vspace{0.8em}
\textbf{Interpretation (parallel slopes).}
\[
\begin{aligned}
T_i=1:\quad & Y(1) \;=\; \alpha + \beta \;+\; \bm{\gamma}^\top \bX_i \;+\; \varepsilon_i,\\[3pt]
T_i=0:\quad & Y(0) \;=\; \alpha \;+\; \bm{\gamma}^\top \bX_i \;+\; \varepsilon_i.
\end{aligned}
\]

\pause
\vspace{0.6em}
\emph{Key assumption:} a common slope \(\bm{\gamma}\) for both groups.
\end{frame}

\begin{frame}{Linear Regression with Covariates (Pooled) (2/3)}
\small

\textbf{Normal equations (univariate \(X\)).} OLS residuals are orthogonal to each regressor:
\[
\sum_{i=1}^n \hat\varepsilon_i = 0, 
\qquad 
\sum_{i=1}^n T_i\,\hat\varepsilon_i = 0, 
\qquad 
\sum_{i=1}^n X_i\,\hat\varepsilon_i = 0.
\]

\pause
\vspace{0.4em}
From the first two,
\[
\bar Y \;=\; \hat\alpha + \hat\beta\,\frac{n_1}{n} + \hat\gamma\,\bar X,
\qquad
\bar Y_1 \;=\; \hat\alpha + \hat\beta + \hat\gamma\,\bar X_1.
\]
(Equivalently, using \(\sum (1-T_i)\hat\varepsilon_i=0\): \(\bar Y_0=\hat\alpha+\hat\gamma\,\bar X_0\).)

\pause
\vspace{0.4em}
\textbf{Eliminate \(\hat\alpha\):}
\[
\hat\beta
\;=\;
(\bar Y_1 - \bar Y_0)\;-\;\hat\gamma\,(\bar X_1 - \bar X_0) \;=\;
\hat\beta_{\text{DiM}}\;-\;\hat\gamma\,(\bar X_1 - \bar X_0).
\]

\pause
\vspace{0.6em}
\textbf{Bias (finite sample).} With pre-treatment \(X\) and CRD,
\[
\E[\hat\beta]
\;=\; \E\!\big[Y(1)-Y(0)\big]
\;-\; \underbrace{\E\!\big[\hat\gamma\,(\bar X_1-\bar X_0)\big]}_{\neq 0, \Cov(\hat\gamma,\bar X_1-\bar X_0)\neq 0} \neq \E\!\big[Y(1)-Y(0)\big].
\]
Under random sampling and CRD, \(\bar X_1-\bar X_0=O_p(n^{-1/2})\) so the bias vanishes as \(n\to\infty\).
\end{frame}

\begin{frame}{Linear Regression with Covariates (Pooled) (3/3)}


\textbf{Variance (heuristic).}
If \(\E[Y\mid T,X]=\alpha+\beta T+\gamma X\) and homoskedasticity holds,
\[
\Var(\hat\beta_{\text{P}} \mid T, X)
\;\approx\;
\Var(\hat\beta_{\text{DiM}} \mid T, X)\,\big(1 - R^2_{Y\sim X}\big),
\]
where \(R^2_{Y\sim X}\) is from regressing \(Y\) on \((1,X)\).
Thus variance decreases with predictive power of \(X\); without correct specification this reduction is not guaranteed.

\pause
\vspace{1em}
\textbf{Takeaway.}
\begin{itemize}
  \item \(\hat\beta_{\text{P}} \to \beta\) (consistent), but not necessarily unbiased in finite samples.
  \pause
  \item If \(\bX\) is predictive of \(Y\), this regression can reduce variance.
  \pause
  \item However, if the model is misspecified (e.g., treatment and control have different covariate effects), efficiency gains may be lost or even reversed.
\end{itemize}
\end{frame}

\begin{frame}{Linear Regression with Covariates (Interacted) (1/2)}
\small
\textbf{Idea.} The pooled model assumes equal slopes for treatment and control.
To relax this, we allow interaction terms between treatment and covariates.
\pause
\[
Y_i = \alpha + \beta T_i + \bm{\gamma}^\top \tilde{\bX}_i + \bm{\delta}^\top (T_i \tilde{\bX}_i) + \varepsilon_i,
\]
where \(\tilde{\bX}_i = \bX_i - \overline{\bX}\) (demeaned covariates).
\pause

\textbf{Interpretation.}
\begin{align*}
T_i = 1 &:~ Y(1) = \alpha + \beta + \bm{\delta}^\top \tilde{\bX}_i + \varepsilon_i, \\[3pt]
T_i = 0 &:~ Y(0) = \alpha + \bm{\gamma}^\top \tilde{\bX}_i + \varepsilon_i.
\end{align*}
\pause
Treatment and control now have different intercepts and covariate effects.

\pause
\textbf{Imputation form.}
The estimated treatment effect can be written as
\[
\hat{\beta}_{\text{I}}
= \frac{1}{n}\sum_{i=1}^n
\big[T_i(Y_i - \widehat{Y}_i(0)) + (\widehat{Y}_i(1) - Y_i)(1-T_i)\big],
\]
where
\[
\widehat{Y}_i(t)
= \hat\alpha + \hat\beta_{\text{I}}t + \hat{\bm{\gamma}}^\top \tilde{\bX}_i + \hat{\bm{\delta}}^\top (t\tilde{\bX}_i).
\]
\end{frame}

\begin{frame}{Linear Regression with Covariates (Interacted) (2/2)}
\small
\textbf{Assumptions:} SUTVA, complete randomization (CRD), and random sampling.
\pause
\vspace{0.8em}

\textbf{Main results:}
\begin{itemize}
  \item \textbf{Consistency:} 
  \[
  \hat\beta_{\text{I}} \;\xrightarrow{p}\; \tau = \E[Y_i(1) - Y_i(0)].
  \]
  \pause
  \item \textbf{Asymptotic normality:}
  \[
  \sqrt{n}\,(\hat\beta_{\text{I}} - \tau)
  \;\xrightarrow{d}\;
  \mathcal{N}\!\left(0,\,\sigma_{\text{I}}^2\right).
  \]
  \pause
  \item \textbf{Efficiency:}
  \[
  \sigma_{\text{I}}^2 \;\le\; \sigma_{\text{DiM}}^2
  \quad\text{and}\quad
  \sigma_{\text{I}}^2 \;\le\; \sigma_{\text{P}}^2.
  \]
  That is, the interacted regression is at least as efficient as DiM or pooled OLS.
  \pause
  \item \textbf{Variance estimation:}  
  The HC2 (heteroskedasticity-robust) variance estimator is consistent for \(\hat\beta_{\text{I}}\).
\end{itemize}
\pause
\bigskip
\textbf{Key insight:} These results hold \emph{even if the regression model is misspecified}, the interacted (Lin, 2013) approach remains valid under randomization.
\end{frame}


\begin{frame}{Question}
 \centering
  {\Large \textbf{Any Questions?}}\\[1em]
  {\Large \textbf{Good luck!}}\\[1em]
  {\Large You’ve got this! Go show off those causal superpowers.}
\end{frame}






\end{document}
